% Short answer key for Listas/lista-r2c-financial-planning-forecasting.md (15 exercises).
% Compile: pdflatex twice; check: pdffonts (Type 1 only, uni = yes).
% Every number is asserted in lista-r2c-financial-planning-forecasting-solutions_check.py.
\documentclass[11pt,a4paper]{article}
\usepackage[utf8]{inputenc}
\usepackage[T1]{fontenc}
\usepackage{lmodern}   % Type 1 vector fonts (replaces PK bitmaps)
\usepackage{cmap}      % ToUnicode CMap -> copyable text and accents
\usepackage[brazilian,english]{babel}
\usepackage{amsmath,amssymb}
\usepackage[a4paper,margin=2.2cm]{geometry}
\usepackage{enumitem}
\usepackage{fancyhdr}
\usepackage{microtype}

\pagestyle{fancy}\fancyhf{}
\lhead{\footnotesize Finan\c{c}as II -- Corporativas (MPE Insper)}
\rhead{\footnotesize List R2c (financial planning) --- Answer key}
\cfoot{\footnotesize\thepage}
\renewcommand{\headrulewidth}{0.3pt}
\setlength{\headheight}{14pt}
\setlist{topsep=2pt,itemsep=1pt}
\setlength{\parindent}{0pt}
\newcommand{\EFN}{\mathit{EFN}}

\usepackage[hyphens]{url}
\begin{document}
\begin{center}{\Large\bfseries Exercise list R2c --- Financial planning and forecasting}\\[4pt]{\large Answer key (full derivations in \url{lista-r2c-financial-planning-forecasting-solutions.pdf})}\end{center}

\section*{1. COGS by percent of sales (Berk 18.1)}
$c=72{,}000/100{,}000=0.72$; COGS $=0.72\times110{,}000=\$79{,}200$.

\section*{2. Reading EFN (Higgins 3.1, 3.4)}
\textbf{3.1} Negative EFN $=$ projected cash surplus above the minimum; set the plug to zero and add it to
cash (matches p.~405). \textbf{3.4} (a) True: $\partial\EFN/\partial AR=1$ since $\EFN=TA-(L+E)$.
(b) False: same assumptions give the same need. (c) False: an annual forecast is a year-end snapshot;
seasonal needs within the year can still arise.

\section*{3. Corpora\c{c}\~ao Felipe (Ross 4.1, 4.2)}
\textbf{4.1} Sales 26,450, costs 19,205, NI 7,245; assets 18,170 = debt 5,980 + equity 12,190. Plug = dividends
$=7{,}245-(12{,}190-10{,}600)=\$5{,}655$. \textbf{4.2} Retained $3{,}622.5$; equity $14{,}222.5$; debt $5{,}200$;
NAF $=18{,}170-19{,}422.5=-\$1{,}252.5$ (surplus).

\section*{4. Growth rates (Berk 18.16)}
(a) ROA$\times b=50/400=12.5\%$. (b) ROE$\times b=50/250=20\%$. (c) $20\%\times0.6=12\%$.

\section*{5. Universal Products (Higgins 3.5)}
(a) $NFA=2{,}000+500-260=2{,}240$. (b) $E=2{,}950+0.5(156)=3{,}028$. (c) $\EFN=(2{,}750+2{,}240)-(1{,}155+500+3{,}028)=4{,}990-4{,}683=\$307$ thousand. Matches p.~405.

\section*{6. Net new financing (Berk 18.2, 18.3)}
\textbf{18.2} $E_1=300{,}000+45{,}000=345{,}000$, $L_1=130{,}000$; NNF $=500{,}000-475{,}000=\$25{,}000$.
\textbf{18.3} Debt + equity fund $470{,}000$ at $D/E=1/3$: $E=352{,}500$, $D=117{,}500$; issue \$7,500 equity and \$17,500 debt.

\section*{7. Jo\~ao Herdeiro (Ross 4.9--4.11)}
\textbf{4.9} Sales 45,600, costs 22,080, EBT 23,520, tax 7,996.8, NI 15,523.2; dividends 6,240; addition to RE
\$9,283.20. \textbf{4.10} Cash 8.03\%, AR 18.16\%, inventory 20.00\%, CA 46.18\%, fixed 90.79\%, TA 136.97\%,
payables 3.42\%; loans, long-term debt, equity: n/a. \textbf{4.11} NI 14,876.4, addition to RE 8,896.4;
TA $=59{,}857.5$; L\&E $=8{,}295+25{,}000+27{,}846.4=61{,}141.4$; NAF $=-\$1{,}283.90$.

\section*{8. Global Corp. (Berk 18.8--18.11), \$ millions}
\textbf{18.8} (a) 189.108; (b) 1.296; (c) NI $=(11.232-7.7)(1-0.7/2.7)=2.616$; (d) 25.056; (e) 19.98;
(f) 16.524; (g) 122.148; (h) 37.476. \textbf{18.9} $E=22.2+0.5(2.616)=23.508$.
\textbf{18.10} NNF $=183.708-174.184=9.524$. \textbf{18.11} $(1-p)2.616\ge1.832\Rightarrow p\le29.98\%$: payout
falls from 50\% to about 30\%.

\section*{9. KMS vs.\ SGR; new debt; IZAX (Berk 18.17--18.19)}
\textbf{18.17} Faster: sales growth 18.0, 16.8, 16.0, 15.3, 14.8\% vs.\ SGR 3.8, 8.2, 8.3, 8.4, 8.5\% (2011--2015);
book $D/E$ rises from 6\% to 32\%. \textbf{18.18} SGR $=0.12\times0.75=9\%$; new debt $=0.09\times400{,}000=\$36{,}000$.
\textbf{18.19} $g_{int}=0.05\times0.75=3.75\%$; $D/E=100{,}000/315{,}000=0.3175$ (from 0.3333).

\section*{10. Under Armour (Higgins 4.8)}
(a) $g^*=PRA\hat T$: 14.96, 17.41, 20.06, 20.16, 19.95\%. (b) Yes: actual growth (18.1--38.4\%) above $g^*$
every year. (c) Higher margin and, mainly, higher leverage; no dividends. (d) $P=8\%$: 22.8\%;
$\hat T=2.4$: 25.2\%; both: 28.8\% ($>27.1\%$ actual).

\section*{11. Turismo da On\c{c}a (Ross 4.25--4.27)}
\textbf{4.25} NI 136,760, addition to RE 95,732; TA 679,080, L\&E 675,232; NAF $=\$3{,}848$.
\textbf{4.26} Capacity sales $929{,}000/0.8=1{,}161{,}250>1{,}114{,}800$: fixed assets stay 413,000; NAF $=596{,}480-675{,}232=-\$78{,}752$.
\textbf{4.27} $D/E=243{,}000/322{,}900=0.7526$; debt allowed 315,044; NAF $=679{,}080-418{,}632-315{,}044\approx-\$54{,}596$.

\section*{12. R\&E Supplies (Higgins 3.9, 3.8)}
\textbf{3.8} \texttt{D19 = C19 + C19*D4}; \texttt{D37 = C37 + D28}. \textbf{3.9} Solve the circular interest:
$X=K/(1-0.375\times0.1)$. (a) $X=2{,}410.54/0.9625\approx\$2{,}504$ thousand ($\approx\$986$ thousand above
2022's 1,518). (b) COGS 84\%: $\$2{,}243$ thousand ($-10.4\%$). (c) Recession: $\$2{,}959$ thousand
($+18.1\%$; book $+18.2\%$), with the formula's negative dividend; $\approx\$3{,}006$ thousand if dividends are
floored at zero. Matches p.~406.

\section*{13. Deriving the formulas (Ross 4.31--4.33)}
\textbf{4.31} NAF $=Ag-ML(V)b(1+g)=-ML(V)b+(A-ML(V)b)g$. \textbf{4.32} NAF $=0\Rightarrow g=ML(V)b/(A-ML(V)b)=
\text{ROA}\,b/(1-\text{ROA}\,b)$; with debt rising at $D/P$ per retained dollar, $A\to P$ gives
$\text{ROE}\,b/(1-\text{ROE}\,b)$. \textbf{4.33} With $P_e=P_b+b\,NI$:
$\text{ROE}_e b/(1-\text{ROE}_e b)=b\,NI/(P_e-b\,NI)=b\,NI/P_b$; same for ROA when assets grow only by retained earnings.

\section*{14. NAF and growth (Ross 4.28, 4.29)}
\textbf{4.28} NAF$(g)=412{,}815(1+g)-491{,}530$: $-16{,}793$ (15\%), $3{,}848$ (20\%), $24{,}489$ (25\%); zero at
$g=19.07\%$ vs.\ formula 16.16\%, because payables grow with sales (18.78\%) and interest is fixed.
\textbf{4.29} NAF$(g)=416{,}783.86(1+g)-554{,}736.22$: $-54{,}596$ (20\%), $-12{,}917$ (30\%), $7{,}922$ (35\%);
zero at $33.10\%$ vs.\ formula 32.24\%, because interest is fixed (margin rises with sales).

\section*{15. KMS at 0.25\% share growth (Berk 18.12--18.15, 18.20--18.22), \$000}
\textbf{18.12} Production 1,000; 1,076.25; 1,157.63; 1,244.42; 1,337.05; 1,435.84: expansion needed in 2012.
\textbf{18.13} Interest 306, 306, 1,666, 1,666, 1,666; tax shields 107.1, 107.1, 583.1 ($\times3$).
\textbf{18.14} NI 8,038.3; 7,911.2; 8,233.7; 9,557.0; 11,007.4. \textbf{18.15} NWC $=39\%$ of sales; increases
2,907.1; 3,118.9; 3,398.8; 3,704.7; 4,040.8. \textbf{18.20} With capex 5,000 (2011), 25,000 (2012), 8,000 after:
FCF 5,773.1; $-12{,}610.8$; 5,376.9; 6,448.2; 7,610.6. \textbf{18.21} $CV=8.5\times26{,}161.5=222{,}373$.
\textbf{18.22} $V=7{,}995.6+222{,}373/1.1^5+1{,}540.7=7{,}995.6+138{,}076.2+1{,}540.7\approx\$147.6$ million
(vs.\ \$201.0 million in the chapter's 1\% scenario).

\end{document}
